\documentclass[12pt]{article}
\usepackage[T1]{fontenc}
\usepackage[utf8]{inputenc}
\usepackage{lmodern}
\usepackage{textcomp}
\usepackage{enumitem}
\usepackage{geometry}
\geometry{margin=1in}
\begin{document}
\begin{center}
Answer Key - Physics - Undergraduate Level Assessment 3 \
\small (Answers are ordered exactly as in the question paper.)
\end{center}

\section*{Multiple Choice}
\begin{enumerate}[label=\arabic*.]
\item \textbf{Answer:} D
\\ \textit{Options:} A) magnetic flux through S; B) rate of change of the magnetic flux through S; C) time integral of the magnetic flux through S; D) rate of change of the electric flux through S
\\ \textit{Note:} The electric displacement current is defined by Maxwell's equations as a quantity that accounts for the changing electric field in a dielectric medium, which is directly proportional to the rate of change of the electric flux through the surface S.
\item \textbf{Answer:} A
\\ \textit{Options:} A) 4; B) 5; C) 6; D) 20
\\ \textit{Note:} The angular magnification of a refracting telescope is determined by the ratio of the focal lengths of the objective lens and the eyepiece lens, and in this case, with a focal length of 20 cm for the eyepiece and an objective focal length of 80 cm (derived from the separation of the lenses), the magnification is calculated as 80 cm / 20 cm = 4.
\item \textbf{Answer:} C
\\ \textit{Options:} A) There are no changes in the internal energy of the system.; B) The temperature of the system remains constant during the process.; C) The entropy of the system and its environment remains unchanged.; D) The entropy of the system and its environment must increase.
\\ \textit{Note:} In a reversible thermodynamic process, the system and its environment can be returned to their original states without any net change in entropy, ensuring that the total entropy remains constant throughout the process.
\item \textbf{Answer:} B
\\ \textit{Options:} A) n = 3, l = 2; B) n = 3, l = 1; C) n = 3, l = 0; D) n = 2, l = 0
\\ \textit{Note:} The transition from n = 4, l = 1 to n = 3, l = 1 is forbidden because it violates the selection rule that states the change in angular momentum quantum number (Δl) must be ±1 for allowed transitions.
\item \textbf{Answer:} D
\\ \textit{Options:} A) 1; B) 2; C) 3; D) 5
\\ \textit{Note:} The correct answer is 5 because for a given angular momentum quantum number l, the magnetic quantum number m\_l can take on values ranging from -l to +l, resulting in a total of 2 l + 1 allowed values, which for l = 2 is 5.
\end{enumerate}

\section*{True / False}
\begin{enumerate}[label=\arabic*.]
\setcounter{enumi}{5}
\item \textbf{Answer:} True
\\ \textit{Options:} A) True; B) False
\\ \textit{Note:} The resolving power of a grating spectrometer is defined as the ratio of the wavelength to the minimum difference in wavelength that can be resolved, and since a resolving power of 250 indicates that it can distinguish between wavelengths that differ by 2 nm (from 500 nm to 502 nm), the statement is true.
\item \textbf{Answer:} True
\\ \textit{Options:} A) True; B) False
\\ \textit{Note:} The statement is true because the magnetic quantum number m\_l can take on integer values ranging from -l to +l, and for l = 2, this results in the values -2, -1, 0, 1, and 2, giving a total of 5 allowed values.
\item \textbf{Answer:} False
\\ \textit{Options:} A) True; B) False
\\ \textit{Note:} Cooper pairs in superconductors are primarily formed due to attractive interactions mediated by lattice vibrations (phonons) rather than the strong nuclear force, which operates at a much smaller scale and is not relevant to the formation of superconducting pairs.
\item \textbf{Answer:} True
\\ \textit{Options:} A) True; B) False
\\ \textit{Note:} The Hall coefficient is directly related to the type of charge carriers in a semiconductor, as it is positive for p-type (holes as carriers) and negative for n-type (electrons as carriers), thus indicating their sign.
\item \textbf{Answer:} True
\\ \textit{Options:} A) True; B) False
\\ \textit{Note:} This statement is true because the energy radiated per second per unit area by a blackbody is proportional to the fourth power of its absolute temperature, as described by the Stefan-Boltzmann law, so increasing the temperature by a factor of 3 results in an increase in energy by a factor of $3^4$, which equals 81.
\end{enumerate}

\section*{Long Form Response}
\begin{enumerate}[label=\arabic*.]
\setcounter{enumi}{10}
\item \textbf{Answer:} To determine the final voltage across a charged capacitor when it is connected to two uncharged capacitors in series, we start by considering a charged capacitor with an initial voltage of V\_0. When this capacitor is connected to two uncharged capacitors in series, the total capacitance of the system can be calculated using the formula for capacitors in series: 1/C\_total = 1/C\_1 + 1/C\_2, where C\_1 and C\_2 are the capacitances of the two uncharged capacitors. The charge Q on the charged capacitor remains constant and is equal to C\_0 * V\_0, where C\_0 is the capacitance of the charged capacitor. After connecting the capacitors, the final voltage V across the charged capacitor is determined by V = Q/C\_total. Through these calculations, it can be shown that the final voltage across the charged capacitor is 2 V\_0/3, reflecting the redistribution of charge across the capacitors in the circuit.
\\ \textit{Note:} correct because it accurately describes the process of calculating the total capacitance of the system when a charged capacitor is connected to uncharged capacitors in series, which is essential for determining the final voltage across the charged capacitor using the principles of charge conservation and the relationship between charge, capacitance, and voltage.
\item \textbf{Answer:} In special relativity, the principle of length contraction states that an object in motion will appear shorter along the direction of its motion to a stationary observer. The formula for length contraction is given by \( L = L_0 \sqrt{1 - \frac{v^2}{c^2}} \), where \( L_0 \) is the proper length, \( v \) is the relative velocity, and \( c \) is the speed of light. To find the velocity required for an observer to measure a 1.00 m rod as 0.80 m, we set up the equation \( 0.80 = 1.00 \sqrt{1 - \frac{v^2}{c^2}} \). Solving for \( v \) leads us to the conclusion that the required velocity is approximately \( 0.60 c \). This demonstrates how relativistic effects can significantly alter measurements depending on the relative motion between observer and object.
\\ \textit{Note:} The answer correctly applies the principle of length contraction from special relativity, using the appropriate formula to relate the proper length of the rod to its contracted length as observed by a stationary observer, while also setting up the equation to solve for the required velocity.
\end{enumerate}

\vfill
\noindent\textit{End of Answer Key}
\end{document}